\documentclass[landscape,fullpage]{article}

\parindent=-2em

\usepackage{eclbkbox}
\usepackage{amsmath}
\usepackage{amssymb}
\usepackage{amscd}
\usepackage{xy}
\usepackage{graphicx}
\usepackage{fancyhdr}
\usepackage{color}
\usepackage[dark,all,bottom,landscape,timestamp]{draftcopy}
% \usepackage{everypage}

\usepackage{ulem}
% go back to italics for emphasis, though
\normalem


% \AddEverypageHook{\waterb}

\fancyhead[LE,RO]{\slshape page number}
\fancyhead[RE]{\slshape \LaTeX}
\fancyhead[LO]{\slshape \LaTeX}
\fancyfoot[LO,LE]{\slshape ! LaTeX Error: Something's wrong---perhaps a missing $\backslash$item?}
\fancyfoot[C]{}
\fancyfoot[RO,RE]{\slshape $\heartsuit\hspace{-1.5em}\longrightarrow$ }

\draftcopyName{DRAFT!}{300}
\pagestyle{fancy}

\begin{document} 


\title{ {\bf How to teach the kids logical frameworks with arcade VideoSystem } }
\author{Tom~Murphy~VII \and George Frankly \and Kate Monday\footnote{Names changed to protect the innocent}}
\date{The year of our lord \today}
\thanks{but no thanks} This page is unintentionally blank but I can't figure out why or what is going on
% \subjclass{subj class}
% \address{127.0.0.1}

% copyright notice generated automatically from Year and author.
% permission added if \permission{} given.

\maketitle

\begin{tabular}{c|c}
\rotatebox{18}{\bf \Large Negative feedback helps the kids to learn.}&
\rotatebox{18}{\bf \Large Negative feedback helps the kids to learn.} 
\end{tabular}


\section{\Large Introduction}

\raggedright

Fact: *LF is dying. In order to attract more students to the study of
logical frameworks, we must approach to them at younger and younger ages.
This work, which funded by the Public Broadcasting Service and Canadian
Illuminati television production company YTV, seeks to capture the
highly sought-after 6--11 year old demographic.

The pitch is simple: A video game so enthralling, that the kids while
their lives away playing it in the videos arcade, all that time
conditioning their mind or minds to reflexively accept dependent
typing; Twelf concrete syntax; the pattern fragment; to use
destination passing style in their everyday language, e.g. to the
mother or mothers: ``Mother, I want {\sf \Large X}. {\tt eq} {\sf
\Large X} {\tt candy};'' to literally {\em not see} the names of bound
variable or variables, just as I can often not remember the eye-color
of the eye or eyes in a girlfriend's face eye holes; to submit to
unification into {\tt \%total}itarian regimes; to only find true
satisfaction in the pale amber glow of a terminal reading {\tt \Large
\%\% SERVER OK \%\%}. Imagine such a world!
%
But this is no pitch, no hypothesis, no hyperbole, no hyperkind. This
game is real or real-seeming. The future is soon or {\em now}.

\vspace{3em}
\verb+\section{LFMan, a video arcade for teaching dependent lambda calculus}+
\vspace{1em} 

LFMan is a game for a single ripe impressionable player or players.
The player plays the game, with the object being to win. A
screen-picture is in
Figure~\verb+Citation 'screenshot' on page 2 undefined on input line 80+.
When the user~\cite{users} or users presses the joystick, the LF Man
moves in the maze [fig.\ref{fig:screenshot_all}] to find dots. The maze is trivial to solve. The
screen and sound presents an extra challenge: It implores the LF Man,
controlled by the player or players~\cite{players}, to evaluate the well-typedness of
LF signatures expressed in Twelf notation strewn roughshod throughout
the maze. A new LF signature awaits the child or childs at every turn.
The child or children must evaluate the signature with
lightnings-quick reflex, to adjust the emacs status bar adorning the
LF Man's banner message. The bar can be as in Figure~\ref{fig:man_examples}.

\begin{figure}[h!t]
\includegraphics[width=0.75 \textwidth]{screenshot}
\caption{The LF Man is faced with a signature. How shall his emacs react?} \label{fig:screenshot}
\end{figure}


% \begin{tabular}{cc}
% \includegraphics[width=2in]{lfman_ok} &
% \includegraphics[width=2in]{lfman_abort}
% \end{tabular}

\begin{figure}[h!t]
\includegraphics[width=\textwidth]{man_examples}
\caption{The bar can be as in Figure~\ref{fig:man_examples}.} \label{fig:man_examples}
\end{figure}


The user or users selects the status bar with the key that has been
painted with the {\sf \large 1} or {\sf \large 2} digits. If there are
2 digits then this is {\sf ABO{\sc R}TS}. The LF Man drives
additionally with the arrows ($downarrow$, $\longleftarrow$,
$\Rightarrow$) painted buttons.

When the LF Man comes lightning-quick to a LF signature, (Figure~~\ref{fig:screenshot})his emacs bar
must accord with the type-correctness of the LF signature. I.E. if the
signature is ill-type, then his emacs bar says ``{\sf ABORTS''} but if
the sigature is well-type, then his emacs {\tt bar says `` OK SERVER
''.} 

\begin{figure}[h!t]
\includegraphics[width=0.75 \textwidth]{screenshot_all}
\caption{This is a picture of the phosphor in the video screen.} \label{fig:screenshot_all}
\end{figure}

When the player's bar accords, the screen and sounds rewards his or
her lightning reflex or reflexes with a Positive Feedback. When
discord arises, then there is Negative Feedback.


\section{Negative Feedback}

\Large \centering When there is Negative Feedback, the mascot of the
Coq theorem prover comes to get the LF Man . This is a chicken. Most
of the time the chicken is spinning and grumbling about H.O.A.S.. But
when there is Negative Feedback the chicken comes to get the LF Man.
If he touches the LF Man then LF Man starts spinning and fading into
the maze or mazes.
The~LF~Man~can~keep~running~and~standing~on~Signatures~or~eating~dots
but he cannot exorcise the chicken

{\it \LARGE
Negative feedback helps the kids to learn.}

\section{Implementation and acknowledgements}

\raggedleft \normalsize The LF Man game is implemented in ActionScript
with Flash VideoSystem.~~~~~~~The source code or codes is shown in
({\sc Figure~\ref{fig:sourcecode}}.) 

The authors are particularly grateful to the Internet People for their
help in learning ActionScript semantics. The best resource was
``steves tutes'', which can be found at
\verb+http://video-animation.com/flash_16.php+. Here 
we learned~~~about inheritance:

\begin{quotation}
When you use Inheritance you save time , because the base attributes and methods already exists,
And Reduce errors , because the base class , Shape has already been used and tested.
And You already know how the base class works. 
Using the same method name to implement different code in subclasses is called PolyMorphism. 

polymorphism does work. .. Try out all the other methods . Try and bust it.. bend it.. shake it..
\end{quotation}

We used inheritance in our VideoSystem application, where the Chicken
class is a subclass of the Man class, and the LFMan class is also a
subclass of the Man class, and the Man class has the implementation of
MovingInTheMazeOrMazes and other shared functionals.
Unfortunately:~~~we observed repeated degenerate behavior where the
Chicken behaves as an LF Man, mimicking his movements and warp-motion
to his locale. We tried using code from {\em steves tutes} to help us

\begin{quotation}
Here is a clas my good friend Lithium wrote as an example.
$\qquad\vdots$ \\
\verb+  var poop = _root.createEmptyMovieClip(name, level);+ \\
\verb+  poop.lineStyle(2, 0x000000, 100);+ \\
$\qquad\vdots$
\end{quotation}

Also unfortunately, this did not accord. In addition to introducing
compile error messages which we do not understand (but that do not
rampart our program from running), it created line-drawings on the
screen that did not appear desirable or related to the chicken
mimicking problem.

Instead we developed a new technique of inheritance which we call {\em
includeitance}. Instead of using the \verb+extends+ keyword to create
a subclass or subclasses, we use the \verb+#include+ keyword at the
head of the class so that LFMan and Chicken can have the same maze
motion but the Chicken does not think he is also a LFMan.

\begin{figure}[t]
\includegraphics[width=\textwidth]{sourcecode}
\caption{All source code. Source is \copyright\ 2009 } \label{fig:sourcecode}
\end{figure}

\begin{breakbox}
Did you want to explore a maze with a chicken? If you have access to children and the children have access to web
browser, put them to this URL: {\bf \LARGE http://spacebar.org/lfman/}
\end{breakbox}

\section{\Large Conclusion\ldots\ or is it?}

NetCraft confirms it. *LF is dying.
I believe children are our future or futures.

\end{document}